\providecommand{\kcat}{\kappa_{\mathrm{cat}}}
\providecommand{\kch}{\kappa_{\mathrm{ch}}}
\providecommand{\kBKM}{\kappa_{\mathrm{BKM}}}
\providecommand{\kfiber}{\kappa_{\mathrm{fiber}}}
\providecommand{\LII}{\Lambda^{2,1}_{\mathrm{II}}}
\providecommand{\PII}{\mathcal{P}_{\mathrm{II}}}
\providecommand{\CII}{\mathcal{C}(\LII)_+}
\providecommand{\DII}{\Delta(\LII)_+}
\providecommand{\Dfive}{D_5}
\providecommand{\DeltaFive}{\Delta_{5}}
\providecommand{\gBKM}{\mathfrak{g}_{\Delta_{5}}}
\providecommand{\Wtwo}{W^{(2)}(\LII)}
\providecommand{\phiEZ}{\phi_{0,1}^{\mathrm{EZ}}}
\providecommand{\ZKthree}{Z_{\mathrm{K3}}}
\providecommand{\CoHA}{\operatorname{CoHA}}
\providecommand{\Yplus}{Y^+}
\providecommand{\Wone}{\mathcal{W}_{1+\infty}}
\providecommand{\Cthree}{\mathbb C^3}

\section*{The \texorpdfstring{$\Delta_{5}$}{Delta5} Denominator and Its BKM Constants}

\paragraph{Theorem: denominator normalization.}
The monic denominator form is
\[
  \Dfive(Z)=64^{-1}\DeltaFive(Z).
\]
The denominator of the Gritsenko--Nikulin generalized
Borcherds--Kac--Moody superalgebra is
\[
  \operatorname{den}(\gBKM)=\Dfive(2Z)=64^{-1}\DeltaFive(2Z).
\]
The primitive denominator has weight \(5\).  The scalar Igusa square has
weight \(10\) and equals the square of the monic Borcherds product.

\paragraph{Fourier input.}
The root exponents are read from the Eichler--Zagier generator
\[
  \phiEZ(\tau,z)=\sum_{n,l} f(n,l)q^n r^l
  =(r^{-1}+10+r)
  +q(10r^{-2}-64r^{-1}+108-64r+10r^2)+O(q^2).
\]
Thus
\[
  c(-1)=1,\qquad c(0)=10,\qquad c(3)=-64,\qquad c(4)=108,
\]
where \(c(4n-l^2)=f(n,l)\).  For
\(\alpha=\alpha(n,l,m)\) in the active positive cone,
\[
  \operatorname{sdim}\gBKM_{\alpha}=f(nm,l)=c(4nm-l^2).
\]
The K3 elliptic genus is
\[
  \ZKthree(\tau,z)=2\,\phiEZ(\tau,z).
\]
The primitive denominator root exponents are the coefficients of
\(\phiEZ\).  The doubled K3 elliptic genus gives the square of the
monic product:
\[
  \operatorname{Borch}(2\,\phiEZ)=\Dfive^2=4096^{-1}\DeltaFive^2.
\]

\paragraph{Weyl sum.}
For \(n,l,m\in\mathbb Z\) with \(n,m>0\),
\[
  a=(n-1)f_2-\frac{l-1}{2}f_3+(m-1)f_{-2}
  \in(\LII)^*=\frac12\LII,
  \qquad
  m(a)=-\frac1{64}f_{\Delta_{5}}(n,l,m),
\]
where \(f_{\Delta_{5}}(n,l,m)\) denotes the corresponding
\(\DeltaFive\)-Fourier coefficient.  On the cone selected by
\[
  \DII^*\subset \CII\subset \DII^{**},
  \qquad
  \mathbb R_{\ge0}\CII^*\cap\CII=\mathbb R_{\ge0}\PII,
\]
the support reduces to \(a\in\LII\cap\mathbb R_{>0}\PII\), with
\(a=0\) contributing \(m(0)=-1\).  The denominator identity is
\[
  \Dfive(z)=
  \sum_{w\in\Wtwo}\det(w)
  \left(
    e^{-\pi i(w\rho,z)}
    -
    \sum_{a\in\LII\cap\mathbb R_{>0}\PII}
    m(a)e^{-\pi i(w(\rho+a),z)}
  \right).
\]

\paragraph{Isotropic ray.}
The three primitive isotropic vertices
\[
  2f_2,\qquad 2f_{-2},\qquad 2f_2-2f_3+2f_{-2}
\]
are permuted by \(\operatorname{Aut}(\PII)\).  For
\(a_0=2f_2\),
\[
  1+\frac1{64}\sum_{t\ge1}f_{\Delta_{5}}(1+2t,1,1)q^t
  =
  \prod_{k\ge1}(1-q^k)^9,
\]
and the isotropic-root coefficients are determined by
\[
  1-\sum_{t\ge1}m(ta_0)q^t
  =
  \prod_{t\ge1}(1-q^t)^{\tau(ta_0)}.
\]
The primitive coefficient is
\[
  \tau(a_0)=9=c(0)-c(-1)=10-1.
\]
Equivalently, the first nontrivial \(\DeltaFive\)-coefficient on this ray
satisfies
\[
  \frac1{64}f_{\Delta_{5}}(3,1,1)=-9,
  \qquad
  f_{\Delta_{5}}(3,1,1)=-576.
\]

\paragraph{Borcherds weight.}
Borcherds' weight theorem gives the weight of the multiplicative lift as
one half of the discriminant-zero input coefficient.  For the level-one
denominator,
\[
  \kBKM(\DeltaFive)=\operatorname{wt}(\DeltaFive)
  =\frac{c_1(0)}2=\frac{10}{2}=5.
\]
The primitive CHL denominator row \(N\in\{1,2,3,4,6\}\) is
\[
\begin{array}{c|ccccc}
N&1&2&3&4&6\\ \hline
c_N(0)&10&8&6&4&2\\
\kBKM(\Phi_N)&5&4&3&2&1
\end{array}
\]
with \(\kBKM(\Phi_N)=c_N(0)/2\) in every row.
The Gritsenko--Clery diagonal-divisor catalogue is a different
triple-indexed table.  Its eight weights and single-cusp constants are
\[
  (5,2,3,1,2,\tfrac12,\tfrac32,1),
  \qquad
  (10,4,6,2,4,1,3,2).
\]

\paragraph{Invariant separation for \texorpdfstring{$K3\times E$}{K3 x E}.}
For \(Y=K3\times E\), the denominator constant determines the BKM entry
\[
  \kBKM(\DeltaFive)=5.
\]
The compact total-space invariants and the specialization invariants are
\[
  \kcat(Y)=\chi(\mathcal O_{K3})\chi(\mathcal O_E)=2\cdot0=0,\qquad
  \kch(Y)=\sum_q(-1)^q h^{0,q}(Y)=1-1+1-1=0,
\]
\[
  \kch^{\mathrm{Heis}}(Y)=3,\qquad
  \kfiber(Y)=24.
\]
The value \(0\) occurs twice, once through Kunneth multiplicativity and
once through the compact Hodge supertrace.  The four numerical values
\(\{0,3,5,24\}\) arise from five subscripted constructions.

\paragraph{Scope theorem.}
For \(d=3\), \(\Phi_{3}\) lands in \(E_1\)-chiral algebras.
The \(E_2\)-structure belongs to centers, doubles, module categories,
and stable-envelope enhancements.  On the toric chart,
\[
  \CoHA(\Cthree)=\Yplus(\widehat{\mathfrak{gl}}_1).
\]
The completed double, center, completion, and evaluation layers produce
\(\Wone\).  The Hall--Drinfeld denominator object is \(\gBKM\) together
with its positive half and completion data.  The K3 self-mirror Yangian
branch is the Mukai-current construction with its own generators and
relations.

\paragraph{Proof obligations.}
The denominator data prove the signed root-character formula for
\(\gBKM\), the monic normalization \(D_5=64^{-1}\DeltaFive\), the
isotropic coefficient \(\tau(a_0)=9\), and the BKM weights
\(\kBKM(\Phi_N)=c_N(0)/2\).

A class-\(\mathcal S\) statement requires an independently constructed
curve, defect datum, and comparison functor to the denominator algebra.

A cohomological-Hall statement requires a compact critical CoHA or
quasi-NCCR source, an orientation/Pfaffian square root, and a character
map whose value is \(D_5^2\) or \(D_5^{-2}\) in the specified scalar
convention.

A compact \(K3\times E\) Hall realization requires a finite recognition
theorem identifying the Hall product with the positive half of
\(U(\gBKM)\), the Borcherds--Serre ideal, the Hopf pairing, the
coproduct, and the completed associator.

Adjacent-volume and \texttt{igusa-cusp-form} concordance checks
normalizations and signs.  The construction of the source category,
target algebra, and comparison morphism remains a separate theorem.
